\chapter{Single column tests}

The directory \verb=Testing/scm= holds single column tests of the physics
schemes: a program calls one scheme through its model entry point on a set
of atmospheric columns, without the dynamical core. They serve as
regression tests, run by \code{make check}, and as a quick way to look at
the response of a scheme to its parameters. The Tiedtke cumulus scheme
(\code{icup}~=~5) is the first one covered.

\section{Files}

\begin{itemize}
  \item \verb=scm_tiedtke.F90=: the driver of the Tiedtke scheme;
  \item \verb=scm_tiedtke.py=: the Python wrapper (run, summary, parameter
        sweep, plot, comparison with a reference);
  \item \verb=scm_profiles.txt=: the test columns;
  \item \verb=make_profiles.py=: builds \verb=scm_profiles.txt= from ERA5;
  \item \verb=run_scm_tests.sh=: the test run by \code{make check};
  \item \verb=reference/=: the reference outputs;
  \item \verb=README=: usage notes.
\end{itemize}

\section{The driver}

The columns are laid out as a $1 \times n$ grid: \code{set\_grid} sets all
the index ranges of the domain (\code{jci1:jci2}, \code{jce1ga:jce2ga},
\ldots) to \code{1:n} by \code{1:1}. The driver then sets the model
configuration seen by the scheme, the one of the MOLOCH core
(\code{idynamic}~=~3, \code{ipptls}~=~2, no chemistry), allocates the
cumulus arrays with the model routines (\code{init\_mod\_cumulus},
\code{allocate\_mod\_cu\_tiedtke}), fills a \code{mod\_2\_cum} structure
from the profiles and calls \code{tiedtkedrv} once. The input humidity is
converted from specific humidity to mixing ratio, as the model stores it;
the vertical velocity goes to \code{avg\_ww} and the domain mean pressure
profile to \code{pmean}, as \code{cumulus} does in
\verb=mod_cu_interface.F90=. The outputs are the \verb=cu_*= arrays of
\verb=mod_cu_common=, plus the cloud type \code{ktype} and the updraft
mass flux \code{pmfu}, public in \verb=mod_cu_tiedtke= for this purpose.

The namelist file (default \verb=scm_tiedtke.in=, or the first argument)
has a \code{scmparam} group (\code{infile}, \code{outfile}, \code{dtcum},
\code{dxkm}) and the model groups \code{tiedtkeparam}, complete, and
\code{cldparam}, with only \code{kfac\_deep}, \code{kfac\_shal} and
\code{k2\_const}. The defaults are the ones of \verb=mod_params.F90=: a
change of a default in the model must be repeated in the driver.

The output is plain text: a header with the namelists used, then for each
column a \verb=COLUMN= line (land flag, cloud type, top and base level,
rain and snow rate, column heating and moistening, the ERA5 CAPE and
convective rain of the profile) and one line per level, top to bottom:
height, pressure, updraft mass flux, tendencies of temperature, humidity,
cloud liquid and ice and of the wind components, convective cloud fraction
and detrained condensate.

\section{Test columns}

\verb=scm_profiles.txt= has twelve columns, six per hemisphere, land and
ocean: Sahel, India, Great Plains, West Pacific warm pool, California
stratocumulus, North Atlantic trades; Amazon, Congo, central Australia,
SPCZ, South East Atlantic stratocumulus, South Atlantic trades. They are
built by \verb=make_profiles.py= from ERA5 (September 2023):
\begin{itemize}
  \item the deep convection columns are taken at the 6-hourly time of the
        month with the largest ERA5 convective rain rate at the point, the
        others at a fixed time of a given day;
  \item the vertical grid is the one of the 32 MOLOCH levels of the South
        America test case ($z = a + b\,h_s$), heights above the ground;
  \item temperature, humidity and wind on the ERA5 pressure levels above
        the ground are interpolated in height; below the lowest of them the
        surface layer is well mixed (constant potential temperature,
        humidity and wind), as a model column is after the PBL scheme; the
        pressure is integrated hydrostatically from the ERA5 surface
        pressure; the vertical velocity is the ERA5 monthly mean;
  \item surface fluxes are the ERA5 monthly means, over land scaled by a
        half sine diurnal cycle (peak $\pi$ times the daily mean) at the
        local solar time; the PBL height is diagnosed from the bulk
        Richardson number (critical value 0.25);
  \item the input tendencies are zero, as in the model, where the cumulus
        scheme is the first physics call of the time step.
\end{itemize}
Each column line of the file has the name, position, land flag, date,
surface pressure, surface fluxes, PBL height, the ERA5 CAPE and convective
rain rate, and the orography; then one line per level (height, pressure,
temperature, specific humidity, wind, vertical velocity, input tendencies
of temperature and humidity) and one per half level (height, pressure).
The ERA5 CAPE and rain rate are a reference for the reader, not a target.

\section{The Python wrapper}

\verb=scm_tiedtke.py= writes the namelist, runs the program in a temporary
directory and reads the output. It needs only the Python standard library,
except for the plot (matplotlib). The program is looked for in the
\code{--exe} option, the \code{SCM\_TIEDTKE} environment variable, the
current directory and the directory of the script.
\begin{verbatim}
  scm_tiedtke.py                            defaults, summary table
  scm_tiedtke.py --list                     namelists used
  scm_tiedtke.py --set entrpen_ocn=1.0e-3 --set lmfwstar=.true.
  scm_tiedtke.py --sweep rprc_ocn=1.0e-3,1.4e-3,2.0e-3
  scm_tiedtke.py --plot tiedtke.png         profiles of every column
  scm_tiedtke.py --save run.out --compare reference/scm_tiedtke.out
  scm_tiedtke.py --check a.out b.out        compare two outputs
\end{verbatim}
The parameter names are the ones of the RegCM namelist; the wrapper knows
the group of each one. The summary table lists, per column, the cloud type,
base and top height, maximum mass flux, rain rate, column heating and
moistening.

The comparison of two outputs requires the same cloud type, base and top in
every column, and the real values to agree within a relative tolerance
(\code{--rtol} or \code{SCM\_RTOL}, default $10^{-4}$) plus an absolute one
(\code{--atol} or \code{SCM\_ATOL}, default $10^{-5}$ times the largest
magnitude of the field in the column), to allow for compiler and flag
differences. It stops at the first difference of each field in a column
and returns a non zero exit status if any is found.

\section{Running the tests}

\code{make check} in \verb=Testing/scm= (or in \verb=Testing=, or at the
top level) builds the drivers, linking the model libraries
\verb=Main/cumlib/libcum.a= and \verb=Main/mpplib/libmpp.a= with the
\verb=Share= library, and runs \verb=run_scm_tests.sh=. For each scheme the
script writes a namelist with the default settings, runs the program in
\verb=scm_output/<scheme>= (or \code{SCM\_OUTDIR}) and, if \code{python3}
is available, compares the output with \verb=reference/scm_<scheme>.out=.
The model libraries must have been built first.

After a change of a scheme meant to change its results, the reference is
updated with
\begin{verbatim}
  scm_tiedtke.py --save reference/scm_tiedtke.out
\end{verbatim}
and its \verb=# profiles:= header line set to \verb=scm_profiles.txt=, so
that the reference does not hold a local path.

\section{Adding a scheme}

A test of another scheme needs:
\begin{enumerate}
  \item a driver \verb=scm_<scheme>.F90= reading \verb=scm_profiles.txt=,
        setting the configuration the scheme sees and calling its model
        entry point, with the namelist groups of the scheme;
  \item the program in \code{check\_PROGRAMS} of
        \verb=Testing/scm/Makefile.am=, with the model library of the
        scheme in its \code{LDADD} and module path in \code{CPPFLAGS};
  \item the scheme in the loop of \verb=run_scm_tests.sh=;
  \item a wrapper, or an extension of \verb=scm_tiedtke.py=, and the
        reference output in \verb=reference/=;
  \item the new files in \code{EXTRA\_DIST} and the program in
        \verb=.gitignore=.
\end{enumerate}
New columns are added in \code{COLUMNS} of \verb=make_profiles.py=; the
regenerated profile file changes the results, so all the references must
be updated with it.
